\begin{abstract}
Heavy-tailed self-regularization (HTSR) reads the quality of a trained layer from the power-law exponent $\alpha$ of its weight spectrum, and has been used to set layer-wise learning rates and pruning ratios. We ask whether $\alpha$ can also \emph{certify}, \emph{select} and \emph{tune} small spiking inserts in an existing dense model, using a pretrained CLIP ViT-B/32 as the test bed. We (i) profile $\alpha$ for all 148 matrices of the model; (ii) convert the MLP activation of one block at a time into a rate-coded integrate-and-fire population and relate the resulting zero-shot and retrieval damage to the block's $\alpha$; (iii) fine-tune the converted or added spiking block under a distillation objective and track $\alpha$, firing statistics and downstream metrics jointly; (iv) convert blocks cumulatively in $\alpha$-ordered, depth-ordered and random orders; and (v) compare spiking lookup-memory and MLP adapters on frozen embeddings. Two methodological findings are needed to make any of this work on CLIP: a per-channel threshold (the blocks have massive-activation channels that a pooled threshold clips) and a per-channel bias correction (only 2--7\,\% of MLP pre-activations are positive, so a ReLU-like rate code without the GELU negative lobe destroys 20--25 points of zero-shot accuracy, which the correction reduces to 1--2 points at $T{=}4$). \TODO{headline numbers for E1/E2}. Code, caches and all run logs are released.
\end{abstract}
